\section{Background and Motivation}
\justify
The integration of mobile computing, artificial intelligence (AI), and conventional
medicine signifies a pivotal frontier in contemporary digital health. The World Health
Organization (WHO) says that 88\% of all countries use traditional and complementary
medicine (T\&CM), with more than 170 member states reporting the use of herbal medicines,
acupuncture, yoga, indigenous therapies, and other types of traditional medicine. Ayurveda,
or ``Science of Life,'' has been practiced in India for more than 3,000 years. It takes a
holistic view of health and divides human physiology into three bio-energetic forces, or
\textit{Doshas}: \textit{Vata} (Kinetic Energy), \textit{Pitta} (Transformative Energy),
and \textit{Kapha} (Cohesive Energy).

\textit{Dravyaguna Vigyan}, the study of the properties of medicinal substances, is the most
important part of Ayurvedic pharmacology. Ayurveda classifies plants based on their
\textit{Rasa} (Taste), \textit{Guna} (Quality), \textit{Virya} (Potency), \textit{Vipaka}
(Post-digestive effect), and \textit{Prabhava} (Specific action), while Western pharmacology
looks for active chemical compounds like alkaloids and glycosides.
\textit{Ocimum sanctum} (Tulsi) is not only an antibacterial agent; it is categorized as
possessing a \textit{Katu} (Pungent) \textit{Rasa} and an \textit{Ushna} (Heating)
\textit{Virya}, rendering it suitable for balancing \textit{Kapha} and \textit{Vata}
disorders while potentially exacerbating high \textit{Pitta} conditions.

Even though Ayurveda is very deep, it is going through a crisis in the 21st century.
The ability to identify medicinal plants, which was once passed down through oral traditions
(\textit{Gurukula}), is fading. Urbanization has separated most people from their natural
surroundings. Because there are not enough qualified \textit{Vaidyas} (Ayurvedic doctors),
millions of people cannot get personalized diagnosis. This creates an urgent need for a
``Digital Vaidya''---a system that uses mobile technology to make expert knowledge
accessible to everyone.

\newpage
\section{Problem Statement}
The digitization of Ayurveda introduces specific computational challenges that current
solutions do not adequately resolve:
\begin{enumerate}[label={\textbf{\arabic*.}}, leftmargin=*]
    \item \textbf{Taxonomic Ambiguity}: Numerous distinct species possess identical
    nomenclature. For example, ``Brahmi'' can mean either \textit{Bacopa monnieri} or
    \textit{Centella asiatica}, depending on region. A text-only search is insufficient
    and dangerous.

    \item \textbf{Visual Similarity}: Many beneficial plants resemble harmful ones.
    For instance, the Solanaceae family contains both poisonous plants and edible
    vegetables, requiring accurate computer vision to distinguish them.

    \item \textbf{Contextual Disconnection}: Current plant identification apps (like
    PlantNet) provide Linnaean taxonomy but no Ayurvedic context. Knowing a plant is
    \textit{Tinospora cordifolia} is not useful without knowing its \textit{Guduchi}
    properties and dosage.

    \item \textbf{Generative Hallucinations}: LLMs like GPT-4 can provide Ayurvedic
    advice but are prone to hallucinating citations and properties. Direct visual LLM
    identification is presently unreliable for essential medical safety.
\end{enumerate}

To address these challenges, we present \textbf{AyurSpace}. Our contributions are:
\begin{itemize}
    \item \textbf{Hybrid Neuro-Symbolic Architecture}: A pipeline that separates
    ``Identification'' (specialized CNNs) from ``Reasoning'' (Semantic LLMs), minimizing
    hallucination risks.
    \item \textbf{Digitized Dravyaguna Ontology}: Medicinal plant properties in a
    searchable, object-oriented schema with algorithmic safety checks.
    \item \textbf{Quantitative Tridosha Assessment}: \textit{Prakriti Pariksha}
    formalized into a discrete-math replicable scoring algorithm.
    \item \textbf{Open Source Mobility}: Cross-platform Flutter implementation for
    accessibility on affordable devices common in developing countries.
\end{itemize}

\newpage
\section{Organisation of Report}
\justify
\begin{description}
    \item[\textbf{Chapter 1:}] Introduction --- Background, motivation, and report structure.
    \item[\textbf{Chapter 2:}] Literature Review --- Existing systems in botanical computer
    vision, LLMs in healthcare, and digital Ayurveda.
    \item[\textbf{Chapter 3:}] Requirement Gathering --- Software and hardware requirements.
    \item[\textbf{Chapter 4:}] Plan of Project --- Methodology, Gantt chart, and system
    architecture.
    \item[\textbf{Chapter 5:}] Project Analysis --- Use case, class, activity, and sequence
    diagrams.
    \item[\textbf{Chapter 6:}] Project Design --- Data flow diagrams and flowchart.
    \item[\textbf{Chapter 7:}] Implemented System --- Architecture, state management, and
    source code listings.
    \item[\textbf{Chapter 8:}] Result Analysis --- Performance benchmarks, accuracy metrics,
    and UI snapshots.
    \item[\textbf{Chapter 9:}] Conclusion and Future Scope --- Summary and future development
    directions.
\end{description}
